\documentclass[11pt]{article}
\usepackage[a4paper, margin=1in]{geometry}
\usepackage{amsmath, amssymb, amsthm}
\usepackage{microtype}
\usepackage{parskip}
\usepackage{enumitem}
\usepackage[colorlinks=true, linkcolor=blue!50!black, urlcolor=blue!50!black]{hyperref}

\theoremstyle{definition}
\newtheorem*{defn}{Definition}
\newtheorem*{ex}{Example}
\newtheorem*{rem}{Remark}

\title{From Groups to Modules\\[2pt]
\large A pedagogical tour of five algebraic structures}
\author{}
\date{}

\begin{document}
\maketitle

\section*{Roadmap}

We will build five structures, in this order:
\[
\text{group} \;\longrightarrow\; \text{ring} \;\longrightarrow\; \text{field}
\;\longrightarrow\; \text{vector space} \;\longrightarrow\; \text{module}.
\]
Each arrow adds either an operation, an axiom, or a notion of ``scalar multiplication.''
The last arrow is the only \emph{generalization} in the chain: a module is what you
get when you take a vector space and allow the scalars to come from a ring instead
of a field. That step trips most people up the first time, so we will spend extra
care on it.

\section{Group}

\paragraph{The idea.}
A group captures the bare minimum needed for an operation that can be \emph{undone}.
Think: integers under addition (you can subtract), or the symmetries of a square
(you can rotate back).

\begin{defn}
A \emph{group} $(G,\cdot)$ is a set $G$ with a binary operation
$\cdot : G \times G \to G$ satisfying:
\begin{enumerate}[label=(G\arabic*), leftmargin=3em]
  \item \emph{Associativity:} $(a\cdot b)\cdot c = a\cdot(b\cdot c)$ for all $a,b,c\in G$.
  \item \emph{Identity:} there is an element $e\in G$ with $e\cdot a = a\cdot e = a$ for all $a$.
  \item \emph{Inverses:} for each $a\in G$ there is some $a^{-1}\in G$ with $a\cdot a^{-1} = a^{-1}\cdot a = e$.
\end{enumerate}
If, in addition, $a\cdot b = b\cdot a$ always holds, the group is called \emph{abelian}.
\end{defn}

\begin{ex}
$(\mathbb{Z},+)$ is an abelian group: identity $0$, inverse of $n$ is $-n$.
\end{ex}

\begin{ex}
The six symmetries of an equilateral triangle (three rotations and three reflections)
form a non-abelian group, traditionally written $S_3$.
\end{ex}

\section{Ring}

\paragraph{The idea.}
A group has \emph{one} operation. But integers, polynomials, and matrices have
\emph{two} — addition and multiplication — and the two interact through
distributivity. A ring axiomatizes this two-operation world.

\begin{defn}
A \emph{ring} $(R,+,\cdot)$ is a set with two binary operations such that:
\begin{enumerate}[label=(R\arabic*), leftmargin=3em]
  \item $(R,+)$ is an abelian group (its identity is written $0$).
  \item $\cdot$ is associative.
  \item \emph{Distributivity:} $a\cdot(b+c) = a\cdot b + a\cdot c$ and $(a+b)\cdot c = a\cdot c + b\cdot c$.
\end{enumerate}
We will assume rings have a multiplicative identity $1$ (such rings are called \emph{unital}).
\end{defn}

\begin{ex}
$\mathbb{Z}$ with the usual $+$ and $\times$ is the prototypical ring.
\end{ex}

\begin{ex}
Polynomials $\mathbb{R}[x]$ form a ring. So do $n\times n$ matrices over $\mathbb{R}$
(non-commutative when $n\geq 2$).
\end{ex}

\paragraph{What's missing?}
In a ring you can add, subtract, and multiply — but you cannot generally
\emph{divide}. There is no element $\tfrac{1}{2}$ in $\mathbb{Z}$. The next step fixes this.

\section{Field}

\paragraph{The idea.}
A field is a ring in which division by anything nonzero is allowed.
$\mathbb{Q}$, $\mathbb{R}$, and $\mathbb{C}$ are the canonical examples.

\begin{defn}
A \emph{field} $(F,+,\cdot)$ is a commutative unital ring in which every nonzero
$a\in F$ has a multiplicative inverse $a^{-1}\in F$ (so $a\cdot a^{-1} = 1$).
\end{defn}

\noindent Equivalently: $(F,+)$ \emph{and} $(F\setminus\{0\},\cdot)$ are both abelian
groups, glued by distributivity.

\begin{ex}
$\mathbb{Q}$, $\mathbb{R}$, $\mathbb{C}$ are fields.
\end{ex}

\begin{ex}
For each prime $p$, the integers modulo $p$, written $\mathbb{F}_p$, form a finite
field with $p$ elements.
\end{ex}

\begin{rem}
Why \emph{prime}? In $\mathbb{Z}/6\mathbb{Z}$ the element $2$ has no inverse: there is
no $x$ with $2x \equiv 1 \pmod 6$. Primality is what guarantees every nonzero element
is invertible. So $\mathbb{Z}/6\mathbb{Z}$ is a ring but not a field — keep this
example in mind, it will return.
\end{rem}

\section{Vector space}

\paragraph{The idea.}
The setting of linear algebra: things you can add together and \emph{scale by
numbers}. The numbers (``scalars'') come from a field.

\begin{defn}
Let $F$ be a field. A \emph{vector space over $F$} is an abelian group $(V,+)$
together with an operation
\[
F \times V \to V, \qquad (\alpha, v) \mapsto \alpha v,
\]
called \emph{scalar multiplication}, satisfying for all $\alpha,\beta\in F$ and $u,v\in V$:
\begin{enumerate}[label=(V\arabic*), leftmargin=3em]
  \item $\alpha(u+v) = \alpha u + \alpha v$
  \item $(\alpha+\beta) v = \alpha v + \beta v$
  \item $(\alpha\beta) v = \alpha(\beta v)$
  \item $1\cdot v = v$
\end{enumerate}
\end{defn}

\begin{ex}
$\mathbb{R}^n$ with componentwise addition and scaling is the standard vector space
over $\mathbb{R}$.
\end{ex}

\begin{ex}
Continuous functions $C([0,1],\mathbb{R})$ form an (infinite-dimensional) vector
space over $\mathbb{R}$.
\end{ex}

\paragraph{The payoff.}
Because the scalars form a field, every vector space has a \emph{basis}, and any
two bases have the same size (the \emph{dimension}). This is what makes linear
algebra so clean.

\section{Module --- the generalization}

\paragraph{The natural question.}
Look back at the vector space definition. Where did we use the fact that $F$ was a
\emph{field}, and not merely a ring? Answer: nowhere in the four axioms.
We only used ``field-ness'' afterwards, when we wanted bases to exist.

So: what if we keep the four axioms but let the scalars come from a mere ring $R$?
That object is called a \emph{module}.

\begin{defn}
Let $R$ be a unital ring. A \emph{module over $R$} (or an \emph{$R$-module})
is an abelian group $(M,+)$ together with an operation
\[
R \times M \to M, \qquad (r, m) \mapsto rm,
\]
satisfying for all $r,s\in R$ and $m,n\in M$:
\begin{enumerate}[label=(M\arabic*), leftmargin=3em]
  \item $r(m+n) = rm + rn$
  \item $(r+s)m = rm + sm$
  \item $(rs)m = r(sm)$
  \item $1\cdot m = m$
\end{enumerate}
\end{defn}

\noindent This is literally the vector space definition with the letter $F$ replaced
by the letter $R$. Nothing else changes.

\subsection*{Reading the phrase ``module over $R$''}

The phrase ``$M$ is a module over the ring $R$'' just names which ring is supplying
the scalars. The letter $R$ is a placeholder — substitute any specific ring you like.

\begin{center}
\begin{tabular}{ll}
\emph{vector space over $\mathbb{R}$} & scalars are real numbers \\
\emph{$\mathbb{Z}$-module} & scalars are integers \\
\emph{module over $R$} & scalars come from some unspecified ring $R$
\end{tabular}
\end{center}

So a ``$\mathbb{Z}$-module'' is not a different beast from ``a module over a
ring $R$.'' It is the special case where $R = \mathbb{Z}$.

\subsection*{Worked example: every abelian group is a $\mathbb{Z}$-module}

This is the example that makes ``scalars from a ring'' click.

Take \emph{any} abelian group $(A,+)$. We will multiply elements of $A$ by integers
in the only way that makes sense, and check that this turns $A$ into a
$\mathbb{Z}$-module --- automatically, for free, with no extra structure required.

For $n\in\mathbb{Z}$ and $a\in A$, define
\[
n\cdot a \;=\;
\begin{cases}
\underbrace{a + a + \cdots + a}_{n \text{ copies}} & n > 0, \\[4pt]
0 & n = 0, \\[4pt]
\underbrace{(-a) + (-a) + \cdots + (-a)}_{|n| \text{ copies}} & n < 0.
\end{cases}
\]
This is just \emph{repeated addition}. The four module axioms (M1)--(M4) are now
forced — they follow from associativity and commutativity of $+$ in $A$. So
\emph{every abelian group becomes a $\mathbb{Z}$-module, in exactly one natural way}.

\begin{ex}[Concrete walkthrough]
Let $A = \mathbb{Z}/6\mathbb{Z} = \{[0],[1],[2],[3],[4],[5]\}$, an abelian group
under addition modulo $6$. Treat it as a $\mathbb{Z}$-module. Then:
\begin{align*}
3 \cdot [2] &= [2] + [2] + [2] = [6] = [0], \\
(-1) \cdot [2] &= [-2] = [4], \\
7 \cdot [3] &= [21] = [3].
\end{align*}
Here $3$, $-1$, $7$ are the \emph{scalars} (elements of $\mathbb{Z}$); $[2], [3]$
are the \emph{module elements} (elements of $\mathbb{Z}/6\mathbb{Z}$). The integer
$3$ acts on $[2]$ by ``add $[2]$ to itself three times.''
\end{ex}

\begin{rem}[The moral]
``Abelian group'' and ``$\mathbb{Z}$-module'' are the same concept in two languages.
The module language just makes the implicit ``repeated addition'' structure visible
as an explicit scalar action.
\end{rem}

\subsection*{What is lost without a field?}

In $\mathbb{Z}/6\mathbb{Z}$ as a $\mathbb{Z}$-module, we saw $3\cdot [2] = [0]$.
There is no way to ``divide by $3$'' and recover $[2]$, because $3\in\mathbb{Z}$
has no multiplicative inverse in $\mathbb{Z}$.

In a vector space (scalars in a field) this can never happen: if $\alpha v = 0$
with $\alpha \neq 0$, multiply by $\alpha^{-1}$ to conclude $v = 0$.

The consequence: modules need not have bases. The $\mathbb{Z}$-module
$\mathbb{Z}/6\mathbb{Z}$ has no basis at all — every element is killed by $6$.
When a module \emph{does} have a basis, it is called \emph{free}.

\section*{Summary}

\begin{center}
\renewcommand{\arraystretch}{1.25}
\begin{tabular}{lll}
\hline
\textbf{Structure} & \textbf{What it has} & \textbf{Key new feature} \\
\hline
Group        & one operation                  & inverses (can undo) \\
Ring         & two operations $(+,\cdot)$     & distributivity \\
Field        & two operations $(+,\cdot)$     & every nonzero scalar invertible \\
Vector space & $+$ and scaling by a field     & bases exist; dimension defined \\
Module       & $+$ and scaling by a ring      & generalizes vector space; basis may not exist \\
\hline
\end{tabular}
\end{center}

\bigskip

The hierarchy is one of progressively richer \emph{scalars}:
\[
\text{abelian group} \;=\; \mathbb{Z}\text{-module}
\;\;\subset\;\; R\text{-module} \;\;\supset\;\; F\text{-vector space}.
\]
The last $\supset$ is the punchline: vector spaces are the special,
well-behaved case of modules — the case where the scalars happen to form
a field.

\end{document}
